\chapter{Scope and Thesis}

This chapter states the central thesis in a way that can absorb both the mathematical and
engineering sides of the repository.

\begin{definition}[Cohomological typing]
Cohomological typing is the discipline of representing typing information as local
sections over a cover of program observations, transporting that information along
restriction maps, and classifying global success or failure by whether those local
sections glue. In the simplest case the data valued in each site is a type-level fact. In
the richer cases used by \DepPy{}, the data also carries trust labels, provenance,
semantic rubrics, proof obligations, and code artifacts.
\end{definition}

The qualifier ``for absolutely everything'' can now be explained more carefully. It does
not mean there is a single theorem that automatically decides all program-analysis
problems. It means the repository repeatedly discovers the same shape of reasoning:
local observations, overlap compatibility, gluing, obstruction, repair, and trust. That
shape is visible in ordinary type checking, in bug diagnosis, in prompt-driven generation,
in semantic oracles, in equivalence testing, in domain-specific theory packs, and in
zero-change adoption for existing Python code. Once the common shape is visible, the
phrase ``typing'' broadens. It comes to mean ``typed local evidence attached to a site and
moved through a gluing discipline.''

The repository pushes this broadening along two axes. The first axis is expressivity. The
plain type system already includes refinements, dependent function and pair formers,
identity-like objects, indexed families, and subtyping. The second axis is evidence.
Hybrid typing extends the foundational language so the same semantic object can be seen at
the level of intent, structural constraints, semantic judgments, proofs, and code. The
result is a repository in which types are no longer merely formulas over values; they are
multi-layer objects that say what is claimed, how it was checked, how much it should be
trusted, and where its remaining ambiguity lives.

A paper with this scope must also be careful about what it is \emph{not} claiming. It is
not claiming that every runtime execution in the repository passes through one uniform
heavy symbolic cohomology engine. It is not claiming that oracle-mediated semantic
judgments have the same epistemic standing as Lean proofs. It is not claiming that every
module is equally mature. Instead, the claim is architectural: the repository is most
coherent when those workflows are viewed as different instances of one local-to-global,
site-aware, evidence-aware semantics.

\paragraph{The central thesis.}
The central thesis of this monograph is that the repository implements, in partial but
already meaningful form, the following picture:
\[
  \Sem : \Site(P)^{\op} \to \mathsf{RefinementLattice}
\]
for the foundational plain system, and
\[
  \HybridTy : (\Site(P) \times \Layer)^{\op} \to \mathsf{RefinementLattice} \times \mathsf{TrustLattice}
\]
for the hybrid system. The former provides the object language of obligations; the latter
provides the layered evidence model in which those obligations are planned, checked,
rendered, translated, discharged, and assembled.

\paragraph{Three design commitments.}
The thesis comes with three commitments that are visible throughout the code base.
\begin{enumerate}[leftmargin=*, itemsep=0.5em]
  \item \textbf{Native Python first.} The mixed-mode surface is designed so users can stay
  inside ordinary Python syntax and still benefit from rich semantics. Surface forms
  degrade gracefully before compilation.
  \item \textbf{Structured evidence, not silent collapse.} Unknown or heuristic facts should
  remain visible as residual or lower-trust data. Trust levels and provenance are meant to
  prevent false certainty.
  \item \textbf{Local-to-global semantics everywhere.} Bugs, equivalence, hallucination,
  library reasoning, and module planning are all interpreted as local facts that may or
  may not glue.
\end{enumerate}

In short: the scope is large, but the architectural thesis is clean.
